% standard beamer lecture template for slides
% by Derek Huang
\documentclass{beamer}
\usepackage[utf8]{inputenc}
\usepackage{algorithm2e, amsmath, amssymb, amsfonts, graphicx}
% allow section.equation numbering
\numberwithin{equation}{section}
% use boadilla theme
\usetheme{Boadilla}
% remove navigation symbols
\usenavigationsymbolstemplate{}
% get numbered figure captions
\setbeamertemplate{caption}[numbered]
% changes itemize to circle + other things
\useoutertheme{split}
\useinnertheme{circles}

% title page stuff. brackets content displayed in footer bar
\title[Math for Data Science, Part 2]{Math for Data Science, Part 2}
% metadata. content in brackets is displayed in footer bar
\author[Derek Huang (BAC Advanced Team)]{Derek Huang}
\institute{BAC Advanced Team}
\date{February 15, 2021}

% change "ball" bullet to numbered bullet and section title for section
\setbeamertemplate{section in toc}{\inserttocsectionnumber.~\inserttocsection}
% change ball to gray square (copied from stackoverflow; \par needed for break)
\setbeamertemplate{subsection in toc}{        
    \hspace{1.2em}{\color{gray}\rule[0.3ex]{3pt}{3pt}}~\inserttocsubsection\par}
% use default enumeration scheme
\setbeamertemplate{enumerate items}[default]
% required line that fixes the problem of \mathbf, \bf not working in beamer
% for later (post-2019) TeX Live installations. see the issue on GitHub:
% https://github.com/josephwright/beamer/issues/630
\DeclareFontShape{OT1}{cmss}{b}{n}{<->ssub * cmss/bx/n}{}

\begin{document}

% title slide
\begin{frame}
    \titlepage
    \centering
    \includegraphics[scale = 0.1]{../bac_logo1.png}
\end{frame}

% table of contents slide
\begin{frame}{Overview}
    \tableofcontents
\end{frame}

\section{Linear algebra}

\subsection{Matrix essentials}

\begin{frame}{Matrix essentials}
    \begin{itemize}
        \item
        \textit{Definition.} A \textit{matrix} $ \mathbf{A} $ of shape
        $ m \times n $ over a field\footnotemark\footnotetext{
            It turns out that it doesn't matter if $ \mathbb{F} = \mathbb{R} $
            or $ \mathbb{F} = \mathbb{C} $, $ \mathbb{C} $ the complex numbers
            \cite{jacob_linalg}.
        } $ \mathbb{F} $ is a function with domain $ \{(i, j) \in \mathbb{N}^2
        : 1 \le i \le m, 1 \le j \le n\} $ and codomain $ \mathbb{F} $
        \cite{jacob_linalg}. Evaluation at $ (i, j) $ is denoted as $ A_{ij} $
        or $ a_{ij} $ and is the $ i, j $th entry of $ \mathbf{A} $.
        \begin{equation*}
            \mathbf{A} = \begin{bmatrix}
                \ a_{11} & \ldots & a_{1n} \ \\
                \ \vdots & \ddots & \vdots \ \\
                \ a_{m1} & \ldots & a_{mn} \
            \end{bmatrix}
        \end{equation*}

        \item
        \textit{Definition.} A \textit{column vector} $ \mathbf{a} $ with
        $ n $ elements is a $ n \times 1 $ matrix over a field $ \mathbb{F} $
        with elements denoted as $ a_1, \ldots a_n $.
        \begin{equation*}
            \mathbf{a} = \begin{bmatrix}
                \ a_1 \ \\ \ \vdots \ \\ \ a_n \
            \end{bmatrix}
        \end{equation*}
    \end{itemize}
\end{frame}

\begin{frame}{Matrix essentials}
    \begin{itemize}
        \item
        \textit{Definition.} The $ n \times n $ \textit{identity matrix}
        $ \mathbf{I} $ is such that
        \begin{equation*}
            \mathbf{I} \triangleq \begin{bmatrix}
                \ 1 & \ldots & 0 \ \\
                \ \vdots & \ddots & \vdots \ \\
                \ 0 & \ldots & 1 \
            \end{bmatrix}
        \end{equation*}
        If context is not clear, we write $ n \times n $ identity as
        $ \mathbf{I}_n $.

	    \item
        \textit{Definition.} The \textit{transpose} of a matrix
        $ \mathbf{A} \in \mathbb{R}^{m \times n } $ is $ \mathbf{A}^\top
        \in \mathbb{R}^{n \times m} $, where the $ i, j $th element of
        $ \mathbf{A}^\top $ is the $ j, i $th element of $ \mathbf{A} $.
        \begin{equation*}
            \mathbf{A}^\top \triangleq \begin{bmatrix}
                \ a_{11} & \ldots & a_{m1} \ \\
                \ \vdots & \ddots & \vdots \ \\
                \ a_{1n} & \ldots & a_{mn} \
            \end{bmatrix}        
        \end{equation*}

        \item
        \textit{Theorem.} For $ \mathbf{A} \in \mathbb{F}^{m \times p} $,
        $ \mathbf{B} \in \mathbb{F}^{p \times n} $,
        $ (\mathbf{AB})^\top = \mathbf{B}^\top\mathbf{A}^\top $
        \cite{jacob_linalg}.
    \end{itemize}
\end{frame}

\begin{frame}{Matrix essentials}
    \begin{itemize}
        \item
        \textit{Definition.} The \textit{sum} of matrices 
        $ \mathbf{A}_1, \ldots \mathbf{A}_q \in \mathbb{F}^{m \times n }$ is
        such that
        \begin{equation*}
            \sum_{k = 1}^q\mathbf{A}_k = \begin{bmatrix}
                \ \sum_{k = 1}^qa_{k, 11} & \ldots &
                    \sum_{k = 1}^qa_{k, 1n} \ \\
                \ \vdots & \ddots & \vdots \ \\
                \ \sum_{k = 1}^qa_{k, m1} & \ldots & \sum_{k = 1}^qa_{k, mn} \
            \end{bmatrix}
        \end{equation*}
        Here $ a_{k, ij} $ is the $ i, j $-th element of $ \mathbf{A}_k $.

        \item
        \textit{Properties of matrix addition} \cite{jacob_linalg}. For
        $ \mathbf{A}, \mathbf{B}, \mathbf{C} \in \mathbb{F}^{m \times n} $,
        $ a, b \in \mathbb{F} $
        \begin{enumerate}
            \item
            $ \mathbf{A} + \mathbf{B} = \mathbf{B} + \mathbf{A} $
            (commutative).

            \item
            $ (\mathbf{A} + \mathbf{B}) + \mathbf{C} = \mathbf{A} +
            (\mathbf{B} + \mathbf{C}) $ (associative).

            \item
            $ a(\mathbf{A} + \mathbf{B}) = a\mathbf{A} + a\mathbf{B} $ (scalar
            distribution).

            \item
            $ (a + b)\mathbf{A} = a\mathbf{A} + b\mathbf{A} $ (scalar
            distribution).

            \item
            $ a(b\mathbf{A}) = (ab)\mathbf{A} $ (scalar associativity).
        \end{enumerate}
    \end{itemize}
\end{frame}

\begin{frame}{Matrix essentials}
    \begin{itemize}
        \item
        \textit{Definition.} The \textit{product} of $ \mathbf{A} \in
        \mathbb{F}^{m \times d} $, $ \mathbf{B} \in \mathbb{F}^{d \times n} $
        is such that
        \begin{equation*}
            \mathbf{AB} = \begin{bmatrix}
                \ \sum_{k = 1}^da_{1k}b_{k1} & \ldots &
                    \sum_{k = 1}^da_{1k}b_{kn} \ \\
                \ \vdots & \ddots & \vdots \ \\
                \ \sum_{k = 1}^da_{mk}b_{k1} & \ldots &
                    \sum_{k = 1}^da_{mk}b_{kn} \
            \end{bmatrix} \in \mathbb{F}^{m \times n}
        \end{equation*}

        \item
	    \textit{Properties of matrix multiplication} \cite{jacob_linalg}. For
	    matrices $ \mathbf{A}, \mathbf{B} \in \mathbb{F}^{m \times d} $,
	    $ \mathbf{C}, \mathbf{D} \in \mathbb{F}^{d \times n} $, and a scalar
	    $ a \in \mathbb{F} $,
	    \begin{enumerate}
	        \item
            $ a(\mathbf{AC}) = (a\mathbf{A})\mathbf{C} $ (scalar
            associativity).

            \item
            $ a(\mathbf{AC}) = \mathbf{A}(a\mathbf{C}) $ (scalar
            commutativity + associativity).

            \item
            $ (\mathbf{A} + \mathbf{B})\mathbf{C} =
            \mathbf{AC} + \mathbf{BC} $ (matrix distribution).

            \item
            If $ \mathbf{I} \in \mathbb{F}^{d \times d} $,
            $ \mathbf{AI} = \mathbf{A} $, $ \mathbf{IC} = \mathbf{C} $
            (identity).
	    \end{enumerate}

        \item
        \textit{Theorem.} For $ \mathbf{A} \in \mathbb{F}^{m \times p} $,
        $ \mathbf{B} \in \mathbb{F}^{p \times q} $, $ \mathbf{C} \in
        \mathbb{F}^{q \times n} $, $ (\mathbf{AB})\mathbf{C} = 
        \mathbf{A}(\mathbf{BC}) $ \cite{jacob_linalg}.
    \end{itemize}
\end{frame}

\begin{frame}{Matrix essentials}
    \begin{itemize}
        \item
        \textit{Definition.} Given matrices $ \mathbf{A}, \mathbf{B} \in
        \mathbb{F}^{m \times n} $, their \textit{Hadamard product} is
        \begin{equation*}
            \mathbf{A} \odot \mathbf{B} \triangleq \begin{bmatrix}
                \ a_{11}b_{11} & \ldots & a_{1n}b_{1n} \ \\
                \ \vdots & \ddots & \vdots \ \\
                \ a_{m1}b_{m1} & \ldots & a_{mn}b_{mn} \
            \end{bmatrix}
        \end{equation*}

        \item
        \textit{Definition.} Given matrices $ \mathbf{A}, \mathbf{B} \in
        \mathbb{F}^{m \times n} $, \textit{Hadamard division} is
        \begin{equation*}
            \mathbf{A} \oslash \mathbf{B} \triangleq \begin{bmatrix}
                \ a_{11} / b_{11} & \ldots & a_{1n} / b_{1n} \ \\
                \ \vdots & \ddots & \vdots \ \\
                \ a_{m1} / b_{m1} & \ldots & a_{mn} / b_{mn} \
            \end{bmatrix}
        \end{equation*}

        \item
        \textit{Definition.} Given matrix $ \mathbf{A} \in
        \mathbb{F}^{m \times n} $, $ b \in \mathbb{R} $, the
        \textit{Hadamard power} is
        \begin{equation*}
            \mathbf{A}^{\odot b} \triangleq \begin{bmatrix}
                \ a_{11}^b & \ldots & a_{1n}^b \ \\
                \ \vdots & \ddots & \vdots \ \\
                \ a_{m1}^n & \ldots & a_{mn}^b \
            \end{bmatrix}
        \end{equation*}
    \end{itemize}
\end{frame}

\begin{frame}{Matrix essentials}
    \begin{itemize}
        \item
        \textit{Definition.} If $ \mathbf{A} \in \mathbb{F}^{n \times n} $ is
        \textit{invertible}, $ \exists! \mathbf{A}^{-1} \in
        \mathbb{F}^{n \times n} $ such that
        \begin{equation*}
            \mathbf{A}^{-1}\mathbf{A} = \mathbf{AA}^{-1} = \mathbf{I}
        \end{equation*}
        The $ \exists! $ notation means ``uniquely exists''.

        \item
        \textit{
            Some properties of invertible matrices\footnotemark\footnotetext{
                We'll discuss rank and eigenvalues later in the slides.                
            }.
        } If $ \mathbf{A} \in
        \mathbb{F}^{n \times n} $ invertible,
        \begin{itemize}
            \item
            $ |\mathbf{A}| \ne 0 $ (nonzero determinant).

            \item
            $ \forall \mathbf{b} \in \mathbb{F}^{n \times 1} $,
            $ \exists! \mathbf{x} \in \mathbb{F}^{n \times 1} $ s.t.
            $ \mathbf{Ax} = \mathbf{b} $ (unique solution to linear system).

            \item
            $ \operatorname{rk}(\mathbf{A}) = n $ (full rank).

            \item
            All eigenvalues of $ \mathbf{A} $ are nonzero.
        \end{itemize}

        \item
        \textit{Theorem.} For invertible $ \mathbf{A}, \mathbf{B} \in
        \mathbb{F}^{n \times n} $, $ (\mathbf{AB})^{-1} =
        \mathbf{B}^{-1}\mathbf{A}^{-1} $ \cite{jacob_linalg}.
    \end{itemize}

    % adjust spacing for footnote
    \bigskip
\end{frame}

\begin{frame}{Matrix essentials}
    \begin{itemize}
        \item
        Matrices and vectors have intuitive interpretations.

        \item
        \textit{Examples.}
        \begin{itemize}
            \item
            \textit{Representing data.} Given data points $ \mathbf{x}_1,
            \ldots \mathbf{x}_N \in \mathbb{R}^d $, we can define the matrix
            $ \mathbf{X} \in \mathbb{R}^{N \times d} $ where each $ k $th row
            is $ \mathbf{x}_k^\top $.

            \item
            \textit{System state.} Suppose $ \mathbf{x} \in \mathbb{R}^n $
            represents the
        \end{itemize}
    \end{itemize}
\end{frame}

\subsection{Vector spaces}

\subsection{Eigenpairs}

\section{Vector calculus}

\subsection{Gradient, Jacobian, Hessian}

\begin{frame}{Gradient, Jacobian, Hessian}
    \begin{itemize}
        \item
        \textit{Definition.} For $ f : \mathbb{R}^n \rightarrow \mathbb{R} $,
        the \textit{gradient} $ \nabla f : \mathbb{R}^n \rightarrow
        \mathbb{R}^n $ is such that
        \begin{equation*}
            \nabla f(\mathbf{x}') \triangleq \begin{bmatrix}
                \ \frac{\partial f}{\partial x_1}(\mathbf{x}') \ \\
                \ \vdots \ \\
                \ \frac{\partial f}{\partial x_n}(\mathbf{x}') \
            \end{bmatrix} \triangleq \left[
                \frac{\partial f}{\partial \mathbf{x}}(\mathbf{x}')
            \right]^\top
        \end{equation*}
        Here $ \mathbf{x}' \in \mathbb{R}^n $ and
        $ \frac{\partial f}{\partial \mathbf{x}} : \mathbb{R}^n \rightarrow
        \mathbb{R}^{1 \times n} $ is the \textit{vector
        derivative}\footnotemark\footnotetext{
            The fact that this is a \textbf{row} vector is important in matrix
            calculus.        
        }.

        \item
        \textit{Definition.} For $ \mathbf{f} : \mathbb{R}^n \rightarrow
        \mathbb{R}^m $, its \textit{Jacobian} is
        \begin{equation*}
            \nabla\mathbf{f}(\mathbf{x}') \triangleq \begin{bmatrix}
                \ \frac{\partial f_1}{\partial x_1}(\mathbf{x}') \ & \ldots &
                    \frac{\partial f_1}{\partial x_n}(\mathbf{x}') \ \\
                \ \vdots & \ddots & \vdots \ \\
                \frac{\partial f_m}{\partial x_1}(\mathbf{x}') \ & \ldots &
                    \frac{\partial f_m}{\partial x_n}(\mathbf{x}') \
            \end{bmatrix} \triangleq \begin{bmatrix}
                \ \frac{\partial f_1}{\partial\mathbf{x}}(\mathbf{x}') \ \\
                \ \vdots \ \\
                \ \frac{\partial f_m}{\partial\mathbf{x}}(\mathbf{x}') \
            \end{bmatrix} \in \mathbb{R}^{m \times n}
        \end{equation*}
        Here $ \mathbf{x}' \in \mathbb{R}^n $, $ \mathbf{f} \triangleq
        [ \ f_1 \ \ldots \ f_ m \ ]^\top $, $ f_k : \mathbb{R}^n
        \rightarrow \mathbb{R} $, $ \forall k \in \{1, \ldots n\} $.
    \end{itemize}
    
    % adjust spacing for footnote
    \medskip
\end{frame}

\begin{frame}{Gradient, Jacobian, Hessian}
    \begin{itemize}
        \item
        Gradient is linear approximation of function change along axes.

        \item
        \textit{Definition.} For $ f : \mathbb{R}^n \rightarrow \mathbb{R} $,
        the \textit{Hessian} $ \nabla^2f : \mathbb{R}^n \rightarrow
        \mathbb{R}^{n \times n} $ is s.t.
        \begin{equation*}
            \nabla^2 f(\mathbf{x}) \triangleq \begin{bmatrix}
                \ \frac{\partial^2 f}{\partial x_1^2}(\mathbf{x}) & \ldots &
                    \frac{\partial^2 f}{\partial x_1x_n}(\mathbf{x}) \ \\
                \ \vdots & \ddots & \vdots \ \\
                \ \frac{\partial^2 f}{\partial x_nx_1}(\mathbf{x}) & \ldots &
                    \frac{\partial^2 f}{\partial x_n^2}(\mathbf{x}) \
            \end{bmatrix}
        \end{equation*}
        $ \forall \mathbf{x}^\top \in \mathbb{R}^n $, $ \nabla^2f(\mathbf{x})
        = \nabla^2f(\mathbf{x})^\top $ ($ \nabla^2f $ symmetric).

        \item
        The Hessian approximates the local curvature of $ f $ at a point.
    \end{itemize}
\end{frame}

\section{Convex optimization}

% BibTeX slide for references. should use either acm or ieeetr style
\begin{frame}{References}
    \bibliographystyle{acm}
    \bibliography{../master_bib}
\end{frame}

\end{document}